\section{Sampling and the geometry of the small values}
\label{sec:sampling}

We now sample the polynomial of Proposition~\ref{prop:auxiliary} along
quadratic paths.  The signed packing makes the leading waves occupy
disjoint arcs of the circle.  Their modulus is consequently independent
of the sampling length, which lets us locate its small values before
constructing the correction.  The sampling argument comes from
\cite[Section~5]{OpenAIMinimal}; the gap description below records the
additional structure needed for the lower bound.

Recall that, for a real vector $Y=(Y_0,\ldots,Y_{N-1})$, we write
\[
 U_Y(t)=\frac1{\sqrt N}\sum_{k=0}^{N-1}Y_k\e(kt).
\]

\begin{proposition}[A structured starting construction]
\label{prop:starting}
Let $0<\delta<1/100$, and put
\[
 K_\delta=\sqrt{\frac{(1+\delta)^3}{1-\delta}},
 \qquad b=\frac18.
\]
There are a smooth even function $B:\T\to[0,K_\delta]$ and, for every
positive integer $N$, real numbers $Y_{N,k}\in[-1,1]$ $(0\le k<N)$
and a smooth function $B_N:\T\to\C$ such that
\begin{gather}
 B_N(-t)=\overline{B_N(t)},\qquad |B_N(t)|=B(t),
 \qquad \norm{U_{Y_N}-B_N}_\infty=O_\delta(N^{-1}),
 \label{eq:starting-approximation}\\
 \int_\T B(t)^2\,dt\ge1-7\delta,
 \qquad
 \frac1N\sum_{k=0}^{N-1}Y_{N,k}^2\ge1-8\delta
 \quad\text{for all sufficiently large }N.
 \label{eq:starting-mass}
\end{gather}
The function $B$ vanishes near $0$ and $1/2$.  Its superlevel set
$\{t\in[0,1/2]:B(t)\ge b\}$ is a nonempty finite union of pairwise
disjoint closed intervals of positive length, contained in $(0,1/2)$.
Let $J_j=[u_j,v_j]$, $1\le j\le M$, be the closures of its complementary
intervals in $[0,1/2]$, and put
\begin{equation}
 \label{eq:starting-gaps}
 l_j=v_j-u_j>0,\qquad
 L=\sum_{j=1}^M l_j,\qquad 0<L\le C_0\delta,
 \qquad C_0=10.
\end{equation}
Thus $B\le b$ on the gaps $J_j$ and $B\ge b$ off their union.

An endpoint $z$ of a gap other than $0$ and $1/2$ is called a
\emph{transition endpoint}.  At each such $z$, $B(z)=b$, and on a
neighborhood $V_z\subset(0,1/2)$ there are a real constant $\alpha_z$
and a real quadratic polynomial $\psi_z$ such that, for every $N$,
\begin{equation}
 \label{eq:starting-endpoint-phase}
 B_N(t)=B(t)\e\bigl(\alpha_z+N\psi_z(t)\bigr)
 \quad(t\in V_z),\qquad
 x_z:=\psi_z'(z)\in(0,1).
\end{equation}
All objects except $Y_N$ and $B_N$ are fixed independently of $N$.
Constants denoted by $O_\delta$ may depend on all the data chosen for
this fixed $\delta$.
\end{proposition}

\subsection{A uniform formula for quadratic sums}

The approximation in \eqref{eq:starting-approximation} must hold at
every angle, including where a stationary point enters or leaves the
support of a cutoff.  Lemma~\ref{lem:quadratic} gives this uniformity;
Poisson summation converts it into the following discrete formula.

\begin{lemma}[Quadratic sampling]
\label{lem:sampling-poisson}
Let $\chi\in C_c^\infty(\R)$, let $x_*\in\R$ and
$\lambda\in\R\setminus\{0\}$, and let $J\subset\R$ be compact.
For $\beta\in J$ and $\ell\in\Z$, put
\[
 x_{\beta,\ell}=x_*+\frac{\ell-\beta}{\lambda},
 \qquad
 \omega_{\beta,\ell}=(\beta-\ell)x_*
                  -\frac{(\beta-\ell)^2}{2\lambda}.
\]
As $N$ tends to infinity through positive integers, uniformly for
$\beta\in J$,
\begin{align}
 &\frac1{\sqrt N}\sum_{k\in\Z}\chi(k/N)
   \e\!\left(\beta k+\frac{N\lambda}{2}(k/N-x_*)^2\right)
 \notag\\
 &\qquad=
 \frac1{\sqrt{|\lambda|}}
 \sum_{\ell\in\Z}\chi(x_{\beta,\ell})
 \e\!\left(\frac{\operatorname{sgn}\lambda}{8}
                      +N\omega_{\beta,\ell}\right)
 +O(N^{-1}).
 \label{eq:sampling-poisson}
\end{align}
The sum on the right has nonzero terms in one fixed finite set of
indices for all $\beta\in J$.  The error constant depends only on
$\chi,x_*,\lambda,J$.
\end{lemma}

\begin{proof}
For $\beta\in J$ and $\ell\in\Z$, define
\[
 \phi_{\beta,\ell}(x)
   =\frac{\lambda}{2}(x-x_*)^2+(\beta-\ell)x.
\]
Poisson summation, applied to a smooth compactly supported function,
gives
\begin{align}
 &\frac1{\sqrt N}\sum_{k\in\Z}\chi(k/N)
   \e\!\left(\beta k+\frac{N\lambda}{2}(k/N-x_*)^2\right)
 \notag\\
 &\hspace{30mm}=
 \sqrt N\sum_{\ell\in\Z}
       \int_\R\chi(x)\e\!\left(N\phi_{\beta,\ell}(x)\right)\,dx.
 \label{eq:sampling-poisson-identity}
\end{align}
The factor $N$ from the change of scale combines with the original
$N^{-1/2}$ to give $\sqrt N$.

Choose a fixed integer $Q$ so large that, for $|\ell|>Q$,
\[
 |\phi'_{\beta,\ell}(x)|
   =|\lambda(x-x_*)+\beta-\ell|\ge |\ell|/2
 \qquad(x\in\operatorname{supp}\chi,\ \beta\in J).
\]
Two integrations by parts have no boundary terms and give
\begin{align*}
 \int\chi\,\e(N\phi_{\beta,\ell})
  =\frac1{(2\pi iN)^2}\int
  \left(
   \frac{\chi''}{(\phi'_{\beta,\ell})^2}
   -\frac{3\lambda\chi'}{(\phi'_{\beta,\ell})^3}
   +\frac{3\lambda^2\chi}{(\phi'_{\beta,\ell})^4}
  \right)\e(N\phi_{\beta,\ell}).
\end{align*}
Each integral is $O(N^{-2}|\ell|^{-2})$.  The entire tail in
\eqref{eq:sampling-poisson-identity} is therefore $O(N^{-3/2})$,
uniformly in $\beta$.  This summable bound is needed before adding
the stationary-phase errors.

For each of the finitely many indices $|\ell|\le Q$, expand
\[
 N\phi_{\beta,\ell}(x)
 =\frac{N\lambda x^2}{2}
   -N(\ell-\beta+\lambda x_*)x+\frac{N\lambda x_*^2}{2}.
\]
Apply Lemma~\ref{lem:quadratic} with $T=N$, curvature $\lambda$, and
$u=N(\ell-\beta+\lambda x_*)$.  Its stationary point is
$x_{\beta,\ell}$, and
$\phi_{\beta,\ell}(x_{\beta,\ell})=\omega_{\beta,\ell}$.
After multiplication by $\sqrt N$, the error is $O(N^{-1})$ uniformly
in $\beta$.  Increase $Q$ if necessary so that
$\chi(x_{\beta,\ell})=0$ whenever $|\ell|>Q$ and $\beta\in J$.
Adding the finite set of leading terms and errors proves the formula.
\end{proof}

\subsection{Disjoint waves and their mass}

\begin{proof}[Proof of Proposition~\ref{prop:starting}]
Fix $\delta$, and take $F$, $\mathcal A=\{\pm a_i:1\le i\le r\}$,
and $v$ from Proposition~\ref{prop:auxiliary}.  Write
\[
 F(y)=\sum_{a\in\mathcal A}c_a\e(a\cdot y),
 \qquad c_a=\widehat F(a),\qquad \lambda_a=a\cdot v.
\]
The widths $w_i=|\lambda_{a_i}|$ satisfy
$2\sum_iw_i<1$.  Lemma~\ref{lem:signed-packing}, with all centers
negated, supplies $H\ge1$ and $\theta_1,\ldots,\theta_H\in\T^m$
such that the closed arcs
\begin{equation}
 \label{eq:sampling-packed-intervals}
 I_{a,h}=-a\cdot\theta_h+
 \left[-\frac{|\lambda_a|}{2H},\frac{|\lambda_a|}{2H}\right]
 \pmod1,
 \qquad a\in\mathcal A,\quad 1\le h\le H,
\end{equation}
are pairwise disjoint.  They have length less than one and avoid
$0$ and $1/2$: an arc containing either point would meet its distinct
reflected arc $I_{-a,h}$ there.

Put $x_h^*=(h-1/2)/H$.  Choose
\begin{equation}
 \label{eq:sampling-cutoffs}
 \chi_h\in C_c^\infty\!\left((h-1)/H,h/H\right),\qquad
 0\le\chi_h\le1,\qquad
 \sum_{h=1}^H\int_\R\chi_h^2\ge1-\delta.
\end{equation}
We require each $\chi_h$ to be positive on a single interval; on that
interval it strictly increases to a plateau at height one and strictly
decreases after the plateau.  The plateau has positive length.  Such cutoffs
exist with arbitrarily short boundary layers in each block, giving
the last inequality in \eqref{eq:sampling-cutoffs}.
Choose real lifts of the $\theta_h$, and set
\begin{equation}
 \label{eq:sampling-coefficients}
 Y_{N,k}=\sum_{h=1}^H\chi_h(k/N)
 F\!\left(k\theta_h+
       \frac N2v(k/N-x_h^*)^2\right),\qquad 0\le k<N.
\end{equation}
At most one cutoff is nonzero for each $k$.  Since $F$ is real and
$\norm F_\infty\le1$, these are real numbers in $[-1,1]$.

For $a\in\mathcal A$, $1\le h\le H$, and $\ell\in\Z$, define, for
$t\in\R$,
\begin{equation}
 \label{eq:sampling-phases}
 \begin{split}
 u_{a,h,\ell}(t)&=t+a\cdot\theta_h-\ell,\\
 x_{a,h,\ell}(t)&=x_h^*-
                      \frac{u_{a,h,\ell}(t)}{\lambda_a},\\
 \psi_{a,h,\ell}(t)&=u_{a,h,\ell}(t)x_h^*
                      -\frac{u_{a,h,\ell}(t)^2}{2\lambda_a}.
 \end{split}
\end{equation}
Define the leading waves by
\begin{equation}
 \label{eq:sampling-leading-waves}
 B_N(t)=\sum_{a,h,\ell}
 \frac{c_a}{\sqrt{|\lambda_a|}}\,
 \chi_h\bigl(x_{a,h,\ell}(t)\bigr)
 \e\!\left(\frac{\operatorname{sgn}\lambda_a}{8}
                  +N\psi_{a,h,\ell}(t)\right).
\end{equation}
The sum is locally finite; shifting $t$ by an integer merely
reindexes $\ell$.  It defines a smooth function on $\T$.

Expand $F$ in \eqref{eq:sampling-coefficients} and apply
Lemma~\ref{lem:sampling-poisson} with
$\beta=t+a\cdot\theta_h$, $\lambda=\lambda_a$, $x_*=x_h^*$,
and $\chi=\chi_h$.  The sum over $k$ extends to $\Z$ because the
cutoffs are supported strictly inside $(0,1)$.  As $t$ ranges over a
fundamental interval, all linear frequencies $\beta$ lie in a fixed
compact interval.  There are only finitely many pairs $(a,h)$, so
\begin{equation}
 \label{eq:sampling-uniform-approximation}
 \norm{U_{Y_N}-B_N}_\infty=O_\delta(N^{-1}).
\end{equation}

For a term of \eqref{eq:sampling-leading-waves} to be nonzero,
$x_{a,h,\ell}(t)$ must lie in the $h$th block.  Thus
\[
 |t+a\cdot\theta_h-\ell|<\frac{|\lambda_a|}{2H},
\]
so $t\pmod1$ lies in $I_{a,h}$.  Pairwise disjointness allows at most
one pair $(a,h)$, and its arc length is less than one, so at most one
integer $\ell$ contributes.  Consequently $|B_N(t)|$ does not depend
on $N$.  It equals the smooth nonnegative function
\begin{equation}
 \label{eq:sampling-modulus}
 B(t)=\sum_{a,h,\ell}
 \frac{|c_a|}{\sqrt{|\lambda_a|}}\,
                 \chi_h\bigl(x_{a,h,\ell}(t)\bigr),
 \qquad 0\le B(t)\le K_\delta.
\end{equation}
It vanishes near $0$ and $1/2$ because all its supporting arcs avoid
those points.  As $F$ is real, $c_{-a}=\overline{c_a}$ and
$\lambda_{-a}=-\lambda_a$.  In
\eqref{eq:sampling-leading-waves}, replacing $(t,a,\ell)$ with
$(-t,-a,-\ell)$ changes $u$ to $-u$, keeps $x$ unchanged, and changes
$\psi$ and $\operatorname{sgn}\lambda_a$ to their negatives.  Hence
$B_N(-t)=\overline{B_N(t)}$, and $B$ is even.

The disjointness also makes the mass calculation exact.  On the arc
indexed by $(a,h)$, the change of variables from $x$ to $t$ has
absolute Jacobian $|\lambda_a|$.  Therefore
\begin{align}
 \int_\T B^2
 &=\sum_{a\in\mathcal A}\sum_{h=1}^H
                  |c_a|^2\int_\R\chi_h(x)^2\,dx
 \notag\\
 &=\norm F_2^2\sum_{h=1}^H\int_\R\chi_h^2
 \ge(1-3\delta)^2(1-\delta)\ge1-7\delta.
 \label{eq:sampling-mass-identity}
\end{align}
If $E_N=\norm{U_{Y_N}-B_N}_\infty$, then
\[
 \left|\int_\T|U_{Y_N}|^2-\int_\T B^2\right|
 \le E_N(2K_\delta+E_N)=O_\delta(N^{-1}).
\]
Parseval now gives the second inequality in \eqref{eq:starting-mass}
for all sufficiently large $N$.  This calculation covers repeated
values of $\lambda_a$, since separation occurs in the angle variable.

\paragraph{The gaps and their endpoint phases.}
It remains to locate the small values of the fixed function $B$.
Each arc supports one bump of height $|c_a|/\sqrt{|\lambda_a|}$.
If this height is at least $b$, the set where that bump is at least
$b$ is a closed interval of positive length.  The positive-length
plateau ensures this even when the height equals $b$.  A lower bump,
including one with $c_a=0$, contributes no such interval.
The supporting arcs are disjoint and avoid $0,1/2$, so the claimed
description of the superlevel set follows.  It is nonempty because
\eqref{eq:sampling-mass-identity} gives $\int B^2>b^2$, and evenness
places a nonempty part in each half-circle.  The closures of the
complementary intervals are thus precisely the gaps $J_j$ in the
statement, and every gap has positive length.

Let $E=\{t\in\T:B(t)<b\}$.  The endpoints of the finitely many
intervals have measure zero, and evenness gives $|E|=2L$.
Using $B\le K_\delta$ and \eqref{eq:sampling-mass-identity},
\[
 1-7\delta\le\int_\T B^2
 \le b^2|E|+K_\delta^2(1-|E|).
\]
Thus
\begin{equation}
 \label{eq:sampling-gap-measure}
 L\le\frac{K_\delta^2-1+7\delta}
              {2(K_\delta^2-b^2)}\le10\delta.
\end{equation}
For the last inequality, $K_\delta^2-1\le5\delta$ when
$\delta<1/100$ and $K_\delta^2-b^2\ge1-1/64$.
Also $L>0$, since $B$ vanishes on neighborhoods of $0$ and $1/2$.

Finally, at a transition endpoint $z$, continuity gives $B(z)=b>0$.
There is a neighborhood of $z$ on which exactly one term of
\eqref{eq:sampling-leading-waves} is nonzero.  For its fixed indices
$(a,h,\ell)$ choose $\alpha_z\in\R$ with
\[
 \e(\alpha_z)=\frac{c_a}{|c_a|}
                 \e\!\left(\frac{\operatorname{sgn}\lambda_a}{8}\right),
 \qquad \psi_z=\psi_{a,h,\ell}.
\]
This yields \eqref{eq:starting-endpoint-phase}.  Differentiating
\eqref{eq:sampling-phases} gives
\[
 \psi_z'(z)=x_{a,h,\ell}(z)\in((h-1)/H,h/H)\subset(0,1),
\]
because the cutoff is positive there.  All data have been fixed
before $N$ varies, and every positive integer $N$ was permitted in
the construction.  This completes the proof.
\end{proof}
